\section{\label{sec:conclusions}Conclusions}

On-shell amplitude techniques have proven to be very effective for computing the renormalization group equations of quantum field theories~\cite{Caron-Huot:2016cwu,EliasMiro:2020tdv,Baratella:2020lzz,Jiang:2020mhe,Bern:2020ikv,Baratella:2020dvw,AccettulliHuber:2021uoa,EliasMiro:2021jgu,Baratella:2022nog,Machado:2022ozb}. 
In particular, the on-shell method~\cite{Caron-Huot:2016cwu} relates anomalous dimensions with unitarity cuts. 
As recently discussed in~\cite{Bresciani:2023jsu}, this method can be easily applied also to describe mixings among operators with different dimensions 
and to capture leading mass effects, which are of paramount importance in several phenomenological studies.

In this work, we have extensively applied the above techniques~\cite{Caron-Huot:2016cwu,Bresciani:2023jsu} to the one-loop renormalization of CP-violating interactions, both above and below the electroweak scale,  of an Axion-Like Particle (ALP) with SM fields, reproducing and extending previous results~\cite{Marciano:2016yhf,DiLuzio:2020oah,Bauer:2020jbp,Chala:2020wvs,DasBakshi:2023lca,Galda:2021hbr,Galda:2023qjx}.

In particular, we  first derived the anomalous dimensions for ALP couplings with fermions, $\phi \bar f f$ and $\phi \bar f i \gamma_5 f$, which require a fermion mass insertion. This allowed us to apply the on-shell method~\cite{Caron-Huot:2016cwu} supplemented by the Higgs low-energy theorem to keep track of leading mass effects while still working in a massless formalism~\cite{Bresciani:2023jsu}. 

Then, we considered the renormalization of ALP couplings to photons and gluons, $\phi FF$ and $\phi GG$, along with their CP counterparts, $\phi F \tilde F$ and $\phi G \tilde G$, recovering the well-known result that they renormalize precisely as $FF$ and $GG$ and, hence, just like their related gauge couplings squared. 
The on-shell method shows this property in a simple and elegant way by just inspecting few tree-level amplitudes.
Moreover, we evaluated the RGEs of operators up to dimension-6 emerging after integrating-out the ALP at one-loop level,  
such as the Weinberg operator $GG\tilde G$ and $GGG$, the (chromo-)magnetic and (chromo-)electric dipole moments, 
i.e.,~$\bar f \sigma\!\cdot\! F f$, $\bar f \sigma\!\cdot\! G f$,
$\bar f \sigma\!\cdot\! F i\gamma_5 f$, and $\bar f \sigma\!\cdot\! G i\gamma_5 f$.

A detailed derivation of the anomalous dimension matrix was carried out both with on-shell and standard techniques, 
aiming to closely compare their virtues and shortcomings.

We found that on-shell methods are computationally advantageous as compared to standard ones, owing to the
significantly lower, as well as less challenging, number of required contributions to be computed.
Moreover, the presence of a large number of non-Abelian gauge bosons generally renders calculations with standard techniques lengthy and computationally expensive: checks for gauge invariance were performed, often involving non-trivial cancellation of different Lorentz and gauge structures.
Remarkably, the on-shell method makes more manifest the connection between the anomalous dimension of operators related by symmetries. 
For instance, the knowledge of the anomalous dimension for the operator $\phi GG$ allowed us to immediately 
infer the one for the CP-dual operator $\phi G\tilde G$: this duality is hidden within the standard approach, where CP-dual operators have different Lorentz structures at the level of Feynman rules, and no similarity in the pattern of cancellations among gauge-dependent terms is present.

Our results were obtained by systematically evaluating all phase-space cut-integrals adopting two different integration techniques, namely by angular integration~\cite{Zwiebel:2011bx}, and via Stokes' theorem~\cite{Mastrolia:2009dr}. The latter, motivated by unitarity, offers the advantage of projecting directly the rational coefficient of the 2-point functions out of the double-cut of one-loop integrals, therefore simplifying the evaluation of anomalous dimensions.  
It would be interesting to extend Stokes' integration method for the evaluation of multi-particle phase-space integrals, and to apply it for the evaluation of anomalous dimensions at higher orders by unitarity cuts, which we plan for future developments.